\section{Conclusion}
\label{sec:conclusion}

We have presented Cognition, a deployable knowledge-graph reasoner for
fact verification that combines three capabilities unavailable in any prior
single system: cassette-native immutable storage, Dempster--Shafer mass
with first-class $m(\Theta)$ as a native abstention signal, and MCTS plus
sheaf GNN multi-hop reasoning over typed infon hypergraphs.  The system
runs entirely on commodity CPU hardware (no GPU required), consumes
approximately 4~GB RAM and 17~MB of retrieval model weights, requires no
pre-defined schema, and produces a calibrated four-part verdict $(m_S, m_R,
m_U, m_\Theta)$ for every claim rather than a single label with a softmax
confidence score.  A Strands-powered conversational Analyst agent makes the
system accessible to non-technical users, citing every source it consults.

Our experimental evaluation on HoVer, AVeriTeC, and SciFact supports all
three directional hypotheses.  H1 results show a 24~pp gap at 4-hop depth
between Cognition+GNN (66.0\%) and flat retrieval (42.0\%), supporting the
hypothesis that MCTS traversal captures multi-hop structure that flat
retrieval cannot; the gap is monotonically increasing with depth (6~pp at
2 hops, 15~pp at 3 hops), confirming that structured traversal rather than
the GNN scorer alone is the source of the advantage.  H2 shows a mean
$m(\Theta)$ of 0.90 on NEI claims for Cognition+GNN versus 0.60 for LLM
zero-shot (Spearman $\rho = 0.95$ vs.\ 0.52), confirming that the DS
vacuous element is an effective and structurally grounded abstention signal
that eliminates 8$\times$ more false-positive endorsements than the LLM
baseline (3\% vs.\ 24\%).  H3 shows Cognition+GNN's ECE of 0.05 versus the
NLI classifier's ECE of 0.12 on SciFact, confirming that DS mass
constructed from explicit evidence sources is better calibrated than a
neural softmax trained on fixed label distributions, even after temperature
scaling.

We plan cross-cassette GNN training on real Wikipedia chain pairs and
multimodal infon extensions (image captions, tabular data) as natural next
steps toward a general-purpose knowledge-graph reasoner.
